\subsection{A sharp all-helper block and an odd-path refinement}

The constant $8/3$ in the proposed all-helper template bound cannot be
increased. We give a thirty-vertex block attaining it, and determine its
unrestricted dissociation number by a separate block argument.

Let $b=6$ and $V=\{0,\ldots,11\}$. Give indices $0,1,4,5,8,9$ colour zero
and the other indices colour one. Define
\begin{align*}
 P={}&\{\{2j,2j+1\}:0\le j<6\},\\
 Q={}&\{\{4j+1,4j+2\},\{4j+3,4j\}:0\le j<3\}.
\end{align*}
Let the owner graph $B$ be the disjoint union of the cycles
\[
 (0,10,1,3,9,2,0),\qquad(4,11,5,7,8,6,4).
\]
Write $H$ for the resulting physical template: helper $i$ is adjacent to
witnesses $x_{2i},x_{2i+1}$, witness $x_v$ is adjacent to owner $c_v$,
and the remaining edges are $Q$ on witnesses and $B$ on owners.

\begin{theorem}\label{thm:sharp-all-helper}
The graph $H$ is connected and bipartite, has minimum degree two and maximum
degree three, and satisfies
\[
 \beta_I(H)=16=8b/3,\qquad \diss(H)=17.
\]
For $r$ disjoint copies these numbers are $16r$ and $17r$, respectively.
Consequently a universal all-helper lower bound $\beta_I\ge11b/4$ is false,
and forcing all helpers can lose an arbitrarily large number of vertices
even with minimum degree two.
\end{theorem}
\begin{proof}
The physical colours satisfy the template requirements. Every witness and
owner has degree three, and every helper has degree two. The three
helper--witness six-cycles are joined by the owner edges, proving connectedness.

All helpers, owner indices
$\{0,3,4,7,8,9,10,11\}$ and witness indices $\{2,6\}$ form a dissociation
set of size sixteen. For the reverse inequality put $A=P\cup Q$; its square
components are $K_j=\{4j,4j+1,4j+2,4j+3\}$, $0\le j<3$.
By Theorem~\ref{thm:all-helper}, a set of size at least seventeen containing
all helpers would give $U$ with $B[V\setminus U]$ dissociation and
$\nu(A[U])\le1$.

In a square, an independent set extends to an opposite pair, and a set
of matching number one extends to a three-set. Thus $U$ can be enlarged to
have three vertices in one square and an opposite pair in each other square.
If its matching number was zero, choose any square to enlarge to a three-set.
Enlarging $U$ preserves dissociation of its owner complement. Call the square
containing three vertices of $U$ charged, and write $S=V\setminus U$.
In an uncharged square $K_0$, $S$ is either $\{0,2\}$
or $\{1,3\}$, both owner edges. Either vertex $9$ or $10$ extends either
edge to an owner three-vertex path. Similarly the two choices in $K_1$ are
the owner edges $\{4,6\},\{5,7\}$; either vertex $8$ or $11$ extends either
edge to such a path. An opposite pair in $K_2$ contains one vertex of each
of $\{9,10\}$ and $\{8,11\}$.

If $K_0$ is charged, the uncharged $K_1,K_2$ already force a three-vertex
owner path. If $K_1$ is charged, use $K_0,K_2$. If $K_2$ is charged, its
single retained vertex extends the edge in $K_0$ or $K_1$ according to which
of these two pairs contains it. Every case contradicts dissociation of
$B[S]$. Hence $\beta_I(H)=16$.

For the unrestricted upper bound divide the physical graph into three
ten-vertex blocks $H_j$: the two helpers paired within $K_j$, the four
witnesses indexed by $K_j$, and their four owners. Each helper--witness
subgraph is a six-cycle. In $H_0,H_1$ the owner edges are the two opposite
witness pairs transported to owners. Within the induced block $H_2$ the
owners are four pendant vertices. An order-four dissociation set in a six-cycle
consists of two
edges and omits two opposite vertices.

We first show that a dissociation set in $H_0$ or $H_1$ has order at most
six, and every six-set contains an entire owner edge. If no owner edge is
fully selected, at most two owners are selected. With at most one owner
the total is at most five. With two owners, a total of six would require
four six-cycle vertices. Each of those four vertices has a selected cycle
neighbour, so the two owners require their witness vertices to be omitted.
Those two omitted witnesses are opposite on the six-cycle and their owners
are an owner edge, a contradiction. If an owner edge is fully selected,
its two witnesses are omitted. Two selected owners give a total at most six.
For three owners the four remaining cycle vertices cannot all be selected:
they induce two edges, and the third owner's witness belongs to one of them.
For four owners all witnesses are omitted and only the two helpers remain.
These cases prove both assertions.

The central block $H_2$ contributes at most seven vertices. Indeed with at
most two owners its contribution is at most six. With at least three owners
it cannot use four cycle vertices, since those would require all their
selected witnesses' owners to be omitted, leaving at most two owners.
Thus it uses at most three cycle vertices, proving the bound seven.
With no owners its contribution is at most four.

Each of $c_9,c_{10}$ has a neighbour on each owner edge of $H_0$:
$c_{10}$ is adjacent to $c_0,c_1$, and $c_9$ to $c_2,c_3$.
Consequently selecting either central owner precludes every six-set of
$H_0$. Similarly selecting either $c_8$ or $c_{11}$ precludes every six-set
of $H_1$. Let $q$ count how many of these two central pairs are met by a
given dissociation set. For $q=0,1,2$ its total order is at most, respectively,
\[
 6+6+4=16,\qquad 5+6+6=17,\qquad 5+5+7=17.
\]
This counts disjoint physical blocks and uses no resource twice.
An explicit seventeen-set consists of helpers $\{0,1,2,3,5\}$,
witness indices $\{0,2,4,6,8,9\}$ and owner indices
$\{1,5,8,9,10,11\}$. Direct inspection gives maximum selected degree one,
so $\diss(H)=17$. Both optima add over disjoint components, proving the
all-$r$ statement. For $r>1$ these examples are disconnected.
\end{proof}

This block differs from the preceding thirty-vertex example, where
$\beta_I=17$ and the loss concerns a further owner-matching restriction.
Here the all-helper problem itself attains $8b/3$; the unrestricted optimum
is larger. Thus it leaves the universal all-helper and unrestricted $8b/3$
questions open.

\paragraph{Using the odd-path surplus.}
Theorem~\ref{thm:owner-matching} also yields a quantitative refinement of its
fixed-extraction corollary. Retain its notation, and let $o$ count odd-order
owner paths in the actual extracted template, which has $b'$ helpers.

\begin{corollary}\label{cor:odd-path-refinement}
Let $G$ be a finite simple bipartite graph of maximum degree at most three,
and let $S$ be a minimal two-star-forming set. Suppose one fixed disjoint
witness extraction of $S$ admits a $K_4$-free choice, has $b'$ helpers and
$o$ odd-order owner paths, and satisfies $o\ge2\gamma b'$ for some real
$0\le\gamma\le1$. Then
\[
 |S|-\diss(G)\le
 \left\lfloor\frac{(5+\gamma)|V(G)|}{51+12\gamma}\right\rfloor.
\]
For $\gamma=0$ this is $5n/51$; for $\gamma=1$ it is $2n/21$.
The statement bounds $\SF-\diss$ when such an extraction exists for a
maximum-cardinality minimal forming set.
\end{corollary}
\begin{proof}
The retained set of size $L$ is disjoint from the template and has no edges
to it. The owner-matching theorem with $z=0$ therefore gives
$Z=3\diss(G)-3L-8b'\ge o/2\ge\gamma b'$.
Use the same nonnegative fixed-extraction variables as above and
$b'=2d-m+\tau-b_3$, where $d=|S|-\diss(G)$. Substituting into the two exact
witness identities gives
\begin{align*}
 (5+\gamma)n-(51+12\gamma)d
 ={}&(Z-\gamma b')+(5+\gamma)(r+q)\\
 &+(12+3\gamma)(\alpha+\pi)+(7+2\gamma)(e-b_3)\\
 &+(2+\gamma)b_3+(24+6\gamma)\tau\ge0.
\end{align*}
Every summand is nonnegative. Division by $51+12\gamma>0$ and integrality
prove the claim. The surplus is assigned once through $Z-\gamma b'$;
it is not added again to either witness exchange.
\end{proof}
